% macros.tex -- shared math macros and theorem environments.
% \input this AFTER amsmath/amssymb/amsthm are loaded (the preambles do so).
% Compatible with pdflatex and tectonic. No shell-escape, no extra packages.

% ---------------------------------------------------------------------------
% Number sets (blackboard bold).
% ---------------------------------------------------------------------------
\providecommand{\R}{\mathbb{R}}
\providecommand{\N}{\mathbb{N}}
\providecommand{\Z}{\mathbb{Z}}
\providecommand{\Q}{\mathbb{Q}}
\providecommand{\C}{\mathbb{C}}
\providecommand{\F}{\mathbb{F}}

% ---------------------------------------------------------------------------
% Paired delimiters (auto-sizing via \left..\right, with a starred *-variant
% for explicit sizing). Requires amsmath (\DeclarePairedDelimiter).
% ---------------------------------------------------------------------------
\DeclarePairedDelimiter{\abs}{\lvert}{\rvert}
\DeclarePairedDelimiter{\norm}{\lVert}{\rVert}
\DeclarePairedDelimiter{\set}{\{}{\}}
\DeclarePairedDelimiter{\ceil}{\lceil}{\rceil}
\DeclarePairedDelimiter{\floor}{\lfloor}{\rfloor}
\DeclarePairedDelimiter{\inner}{\langle}{\rangle}

% ---------------------------------------------------------------------------
% Calculus helpers.
%   \dd            upright differential "d"
%   \dd[n]{x}      n-th differential of x, e.g. \dd{x} -> dx, \dd[2]{t}
%   \eval          evaluation bar: \eval{expr}{a}{b}
% ---------------------------------------------------------------------------
\providecommand{\dd}{}
\renewcommand{\dd}[2][]{\mathop{}\!\mathrm{d}^{#1}#2}
\newcommand{\eval}[3]{\left.#1\right\rvert_{#2}^{#3}}

% Common operators not predefined by amsmath.
\DeclareMathOperator{\diverg}{div}
\DeclareMathOperator{\curl}{curl}
\DeclareMathOperator{\grad}{grad}
\DeclareMathOperator{\rank}{rank}
\DeclareMathOperator{\trace}{tr}
\DeclareMathOperator{\spann}{span}
\DeclareMathOperator*{\argmin}{arg\,min}
\DeclareMathOperator*{\argmax}{arg\,max}

% Convenience shorthands.
\providecommand{\eps}{\varepsilon}
\providecommand{\To}{\longrightarrow}
\newcommand{\deriv}[2]{\frac{\dd{#1}}{\dd{#2}}}
\newcommand{\pderiv}[2]{\frac{\partial #1}{\partial #2}}

% ---------------------------------------------------------------------------
% Theorem environments (amsthm). Shared numbering within sections.
% ---------------------------------------------------------------------------
\theoremstyle{plain}
\newtheorem{theorem}{Theorem}[section]
\newtheorem{lemma}[theorem]{Lemma}
\newtheorem{proposition}[theorem]{Proposition}
\newtheorem{corollary}[theorem]{Corollary}

\theoremstyle{definition}
\newtheorem{definition}[theorem]{Definition}
\newtheorem{example}[theorem]{Example}

\theoremstyle{remark}
\newtheorem*{remark}{Remark}
% (The `proof` environment is provided by amsthm automatically.)
